\clearpage
\subsection{RNF-01 — Soporte bilingüe Español / Inglés}

\renewcommand{\arraystretch}{1.3}

\noindent
\begin{tabular}{|p{5cm}|p{10.5cm}|}
\hline
\rowcolor{secVerde}\multicolumn{2}{|l|}{\color{white}\textbf{1. INFORMACIÓN DEL CLIENTE}} \\
\hline
Nombre de la empresa  & UTEQ \\ \hline
Nombre del Contacto   & Angélica Garduño Bustamante \\ \hline
Fecha de Elaboración  & 01/06/2026 \\ \hline
\end{tabular}

\vspace{0.3cm}

\noindent
\begin{tabular}{|p{5cm}|p{10.5cm}|}
\hline
\rowcolor{secVerde}\multicolumn{2}{|l|}{\color{white}\textbf{2. INFORMACIÓN GENERAL}} \\
\hline
Nombre del Requerimiento     & RNF 01. Soporte bilingüe Español / Inglés \\ \hline
Tipo de Requerimiento        & No funcional — Usabilidad \\ \hline
Fecha requerida de solución  & 25/06/2026 \\ \hline
Prioridad                    & Media \\ \hline
\end{tabular}

\vspace{0.3cm}

\noindent
\begin{tabular}{|p{4.5cm}|p{11cm}|}
\hline
\rowcolor{secVerde}\multicolumn{2}{|l|}{\color{white}\textbf{3. DESCRIPCIÓN DEL REQUERIMIENTO}} \\
\hline
Objetivo del Requerimiento &
Garantizar que la totalidad de las cadenas de texto visibles para el usuario estén disponibles en español e inglés, permitiendo el cambio dinámico del idioma desde la pantalla de perfil. \\
\hline
Beneficio esperado &
Accesibilidad para usuarios anglófonos, experiencia profesional consistente y posibilidad de uso académico internacional del proyecto. \\
\hline
Alcance &
- Extraer las cadenas hardcoded restantes (\texttt{login}, \texttt{register}, \texttt{onboarding}, alertas) hacia los archivos \texttt{locales/es.json} y \texttt{locales/en.json}. \par
- Localizar el formato de fechas según el idioma activo. \par
- Persistir la preferencia de idioma en \texttt{AsyncStorage}. \\
\hline
Ubicación & \url{https://github.com/MarkG777/AbejaNet} \\ \hline
Descripción de Usuarios y Responsables &
\textbf{Usuarios:} \par
\textbf{Usuario (Cliente):} cambia el idioma desde el perfil. \par
\vspace{0.1cm}
\textbf{Responsables:} \par
\textbf{App Móvil:} aplica el idioma seleccionado a través de \texttt{react-i18next}. \par
\textbf{AsyncStorage:} persiste la preferencia entre sesiones. \\
\hline
Descripción funcional &
1. El usuario accede a su perfil y elige el idioma deseado. \par
2. Toda la interfaz se actualiza al instante sin reiniciar la app. \par
3. La preferencia se almacena localmente y se aplica al siguiente arranque. \par
4. Los formatos de fecha y hora se muestran según la convención del idioma activo. \\
\hline
Restricciones &
- Solo se soportan español e inglés en esta iteración. \par
- Las claves nuevas deben mantener consistencia entre ambos archivos JSON. \\
\hline
Precedencias & Ninguna (la base de \texttt{i18next} ya existe en el proyecto). \\ \hline
Documentación requerida & Manual técnico con la convención de claves de internacionalización. \\ \hline
Observaciones & La librería \texttt{react-i18next} ya se encuentra instalada en el proyecto. \\ \hline
\end{tabular}

\vspace{0.3cm}

\noindent
\begin{tabular}{|p{5.1cm}|p{5.1cm}|p{5.1cm}|}
\hline
\rowcolor{secVerde}\multicolumn{3}{|l|}{\color{white}\textbf{4. CONDICIONES DE PRUEBA}} \\
\hline
\rowcolor{cellGris}\textbf{Condición} & \textbf{Validación} & \textbf{Resultado} \\
\hline
El usuario cambia el idioma de español a inglés desde el perfil. &
La interfaz debe traducirse en tiempo real, incluyendo formularios, validaciones y alertas. &
Todos los textos visibles cambian al inglés sin reiniciar la app. \\
\hline
El usuario cierra y vuelve a abrir la app después de haber cambiado el idioma. &
La preferencia debe persistir entre sesiones. &
La app inicia automáticamente con el idioma previamente seleccionado. \\
\hline
\end{tabular}
